\subsection{Logistic Regression}
\label{sec:logistic_regression}

\begin{table*}[htbp]
\caption{Hyperparameter Search Space for Logistic Regression Classifier.}
\label{tab:5}
\label{tab:V}
\label{tab:lr_params}
\centering
\footnotesize
\begin{tabular}{llc}
\toprule
Component / Hyperparameter & Evaluated Values & Count \\
\midrule
Feature Scaler & None, MinMaxScaler & 2 \\
Dimensionality Reduction (PCA) & None, 0.50, 0.75 & 3 \\
Regularization Strength & 0.01, 0.1, 1, 10 & 4 \\
Optimization Solver & 'lbfgs', 'saga' & 2 \\
Maximum Iterations & 500, 1000 & 2 \\
\midrule
\multicolumn{2}{r}{\textbf{Total Grid Combinations}} & \textbf{96 (480 fits)} \\
\bottomrule
\end{tabular}
\end{table*}

The multinomial regression approach (Softmax regression) in LR maps latent feature representations into six intersectional demographic categories through unified linear probability estimation~\cite{ref37}. The Softmax probability function computes the conditional probability of a sample belonging to the target category as defined in \eqref{eq:10}, where $\mathbf{x}$ represents the input feature vector, $y$ is the class label variable, $c$ denotes the target class category, $\mathbf{w}_c$ and $b_c$ are the weight vector and bias of class $c$, respectively, while the constant $K = 6$ indicates the total number of intersectional classes~\cite{ref38}:
\begin{equation}
P(y = c \mid \mathbf{x}) = \frac{e^{\mathbf{w}_c^T \mathbf{x} + b_c}}{\sum_{j=1}^K e^{\mathbf{w}_j^T \mathbf{x} + b_j}}
\label{eq:10}
\end{equation}
The model training process minimizes the cross-entropy loss function by incorporating an $L_2$ regularization penalty to control weight complexity in the high-dimensional representation space~\cite{ref39}. The structured hyperparameter search evaluates variations in regularization strength $C$, optimization solver algorithms ('lbfgs', 'saga'), the iteration limit max\_iter, and the integration of scaling and PCA, yielding 96 evaluation combinations or 480 total fits under the 5-Fold StratifiedGroupKFold CV protocol presented in Table~\ref{tab:5}.
